% ============================================================
%  Sezione 3.4 — Lo scheduler energy-aware
% ============================================================

\section{Lo scheduler energy-aware}
\label{sec:scheduler}

Lo scheduler energy-aware costituisce il componente decisionale centrale
dell'estensione progettata. Il suo compito è selezionare, tra le varianti
registrate per una data funzione logica, quella da eseguire concretamente
in risposta a ciascuna richiesta di invocazione. La selezione avviene
prima dell'assegnazione del container e dell'avvio dell'esecuzione vera e
propria, consentendo di sfruttare le stime energetiche aggiornate presenti
in \textit{etcd} come base per le proprie decisioni.

A differenza degli scheduler tradizionali, che operano esclusivamente a
livello di selezione del nodo di esecuzione, lo scheduler progettato
interviene a granularità di variante: a parità di funzione logica
richiesta, è possibile eseguire implementazioni alternative caratterizzate
da profili energetici e livelli di errore distinti. Questa scelta
progettuale consente di trattare il trade-off tra consumo energetico e
qualità del risultato come una grandezza controllabile a runtime, anziché
come una proprietà fissa dell'implementazione.

Il comportamento dello scheduler è ispirato ai meccanismi di controllo
runtime introdotti da JouleGuard~\cite{jouleguard}, adattati al paradigma
FaaS e al modello a varianti descritto nella Sezione~\ref{sec:variant-model}.
In particolare, il sistema non applica alcuna selezione alternativa in
assenza di un'esplicita autorizzazione da parte dell'utente, garantendo la
piena compatibilità con il comportamento legacy della piattaforma.

La logica decisionale dello scheduler si articola in tre casi distinti,
determinati dai parametri forniti dall'utente al momento dell'invocazione:
il caso base (senza abilitazione della variant selection), il caso di
selezione per errore minimo e il caso di selezione con vincolo
energetico esplicito. Tali casi sono analizzati nelle sezioni seguenti.

% --- FIGURA SUGGERITA -------------------------------------------
% Diagramma di flusso della logica decisionale dello scheduler:
% dall'invocazione in ingresso, biforcazione su AllowApprox,
% ramo error-first, ramo energy-budget, rifiuto HTTP 422.
%
% \begin{figure}[htbp]
%   \centering
%   \includegraphics[width=0.85\textwidth]{figures/scheduler_flow}
%   \caption{Flusso decisionale dello scheduler energy-aware. A seconda
%            del consenso fornito dall'utente e della presenza di un
%            vincolo energetico, il sistema applica politiche di selezione
%            distinte o preserva il comportamento legacy.}
%   \label{fig:scheduler_flow}
% \end{figure}
% ----------------------------------------------------------------


% ------------------------------------------------------------
\subsection{Il consenso esplicito dell'utente}
\label{sec:scheduler:consenso}
% ------------------------------------------------------------

Un principio fondamentale che guida la progettazione dello scheduler è la
separazione netta tra il comportamento di default della piattaforma e la
modalità energy-aware. La selezione di varianti alternative non viene attivata di default dalla piattaforma:
essa richiede sempre un consenso esplicito da parte dell'utente che
formula la richiesta di invocazione.

A tale scopo, l'interfaccia di invocazione è stata estesa con due campi
aggiuntivi: il flag booleano \texttt{AllowApprox} e il parametro
opzionale \texttt{MaxEnergyJoule}. Il primo esprime il consenso
dell'utente all'utilizzo di varianti alternative rispetto a quella
nominale; il secondo, se presente e positivo, specifica un budget
energetico massimo accettabile per l'invocazione, espresso in Joule.

In assenza di autorizzazione esplicita — ovvero quando
\texttt{AllowApprox = false} — la variant selection è completamente disabilitata e la
richiesta viene instradata verso la variante specificata nella chiamata,
secondo il comportamento originale della piattaforma. Questo garantisce la
piena compatibilità retroattiva con le applicazioni preesistenti e con gli
utenti che non desiderano delegare allo scheduler la scelta
dell'implementazione.

Quando invece \texttt{AllowApprox = true}, lo scheduler attiva la variant
selection energy-aware. In questo caso, tutte le varianti registrate per
la funzione logica corrispondente diventano candidate all'esecuzione e la
selezione avviene secondo criteri energetici e di errore,
eventualmente vincolati dal parametro \texttt{MaxEnergyJoule}. Il flusso
di controllo ad alto livello è formalizzato nell'Algoritmo~\ref{alg:submit}.

\begin{algorithm}[htbp]
\caption{\texttt{SubmitRequest} con variant selection opt-in}
\label{alg:submit}
\begin{algorithmic}[1]
\Function{SubmitRequest}{$r$ : \texttt{*function.Request}}
    \If{$r.\texttt{AllowApprox} = \textbf{true}$}
        \State $(\textit{selectedFn},\; \textit{err}) \gets
               \Call{SelectEnergyAwareVariant}{r}$
        \If{$\textit{err} \neq \textbf{nil}$}
            \State \Return $(\textbf{nil},\; \textit{err})$
            \Comment{vincolo energetico non soddisfatto → HTTP 422}
        \EndIf
        \If{$\textit{selectedFn} \neq \textbf{nil}$}
            \State $r.\texttt{Fun} \gets \textit{selectedFn}$
            \Comment{sostituisce la variante richiesta}
        \EndIf
    \EndIf
    \State $\textit{requests} \leftarrow r$
    \Comment{invia al loop legacy: \texttt{p.OnArrival(r)}}
    \State \Return \Call{SchedulingDecision}{$r$}
\EndFunction
\end{algorithmic}
\end{algorithm}

La struttura dell'algoritmo evidenzia come la variant selection si
inserisca come un livello di pre-elaborazione rispetto allo scheduler
originale: il risultato della selezione è la sostituzione della variante
richiesta con quella ottimale individuata dalla funzione
\textsc{SelectEnergyAwareVariant}, dopodiché l'esecuzione prosegue
secondo il percorso ordinario. In caso di errore — tipicamente quando
nessuna variante soddisfa i vincoli imposti — la richiesta viene rifiutata
con un codice di stato HTTP~422 (\textit{Unprocessable Entity}) senza
avviare alcuna computazione.


% ------------------------------------------------------------
\subsection{Selezione per errore minimo}
\label{sec:scheduler:error}
% ------------------------------------------------------------

Quando l'utente abilita la variant selection (\texttt{AllowApprox = true})
senza specificare un budget energetico, oppure con un valore
\texttt{MaxEnergyJoule} nullo o non positivo, lo scheduler opera secondo
la politica di \textit{accuracy-first}. In questo scenario, l'obiettivo è individuare la
variante che offre il risultato di qualità più elevata tra tutte quelle
disponibili, ricorrendo al consumo energetico stimato esclusivamente nel caso di parità di errore minimo tra due o più varianti.

La politica non applica alcun filtraggio energetico preventivo: tutte le
varianti registrate per la funzione logica sono candidate alla selezione,
indipendentemente dal loro profilo energetico. Questo approccio è coerente
con la natura opt-in del meccanismo: l'utente ha autorizzato l'uso di
varianti alternative, ma non ha imposto alcun vincolo sul consumo; la
piattaforma interpreta quindi tale consenso come un interesse prioritario
verso la qualità del risultato.

Per confrontare varianti con modelli di output eterogenei viene definita
una funzione di punteggio di errore $Err(v)$, costruita in modo tale
che valori numericamente inferiori corrispondano a errore minore.
La funzione è implementata nel componente \texttt{ErrorScore}$(v)$
secondo le regole seguenti:

\begin{itemize}
  \item se il modello di output della variante è di tipo \texttt{"error"},
        il punteggio è pari alla stima dell'errore associato:
        $Err(v) = \textit{ErrorEstimate}(v)$;

  \item se il modello è di tipo \texttt{"quality"}, la classe qualitativa
        viene convertita in una scala numerica ordinale:
        $\textit{High} \to 0$, $\textit{Medium} \to 1$, $\textit{Low} \to 2$;

  \item in assenza di un modello di output esplicito, si pone per
        convenzione $Err(v) = 0$, attribuendo l'errore minimo alla
        variante.
\end{itemize}

La selezione della variante ottimale $v^*$ avviene mediante
minimizzazione lessicografica sulla coppia $(Err(v),\, E(v))$, dove $E(v)$
rappresenta la stima energetica per l'invocazione:

\begin{equation}
  v^* = \arg\min_{v \in V}\; \bigl(Err(v),\; E(v)\bigr)
  \label{eq:error_first}
\end{equation}

In altri termini, tra tutte le varianti disponibili viene selezionata
quella con il punteggio di errore minore; in caso di parità, viene
preferita la variante con il consumo energetico stimato più basso. La
stima $E(v)$ è calcolata dalla funzione ausiliaria
\texttt{EnergyEstimate}$(v, \textit{warm})$, che distingue tra esecuzioni
a caldo (\textit{warm start}) e a freddo (\textit{cold start}):

\begin{equation}
  E(v) =
  \begin{cases}
    \textit{InvocationJoule}(v) & \text{se } \textit{isWarm} = \textbf{true} \\
    \textit{ColdStartJoule}(v) + \textit{InvocationJoule}(v)
      & \text{se } \textit{isWarm} = \textbf{false}
  \end{cases}
  \label{eq:energy_estimate}
\end{equation}

La determinazione dello stato warm viene effettuata dalla funzione
\texttt{WarmHint}$(v)$, che verifica l'esistenza di un container già
attivo utilizzabile per la variante $v$. In particolare, se la variante
espone il flag \texttt{ShareContainer = true}, il sistema considera warm
anche i container di altre varianti appartenenti alla stessa funzione
logica e con lo stesso runtime. Questo meccanismo di condivisione consente di
ridurre l'impatto del cold start nelle situazioni in cui una variante
priva di container attivo può beneficiare di un ambiente di esecuzione già
inizializzato da una variante affine.

I valori \textit{InvocationJoule} e \textit{ColdStartJoule} utilizzati
nella stima sono recuperati da \textit{etcd} al momento della selezione e
corrispondono all'ultima stima EMA disponibile per ciascuna variante,
aggiornata asincronamente dal worker di raccolta descritto nella
Sezione~\ref{sec:worker}. Lo scheduler opera quindi sempre sulla stima più
recente disponibile, con un ritardo rispetto al comportamento reale
dell'ordine di grandezza del periodo del worker.


% ------------------------------------------------------------
\subsection{Selezione con vincolo energetico}
\label{sec:scheduler:budget}
% ------------------------------------------------------------

Quando l'utente fornisce un budget energetico esplicito
(\texttt{MaxEnergyJoule} $> 0$), lo scheduler adotta la politica di
selezione con vincolo energetico, denominata \textit{energy-budget aware}.
In questa modalità il consumo energetico stimato diventa un vincolo
esplicito, e la minimizzazione dell'errore avviene all'interno
dell'insieme delle varianti compatibili con tale vincolo.

Il processo di selezione si articola nelle fasi descritte di seguito.

\medskip
\noindent\textbf{Fase 1 — Partizionamento delle varianti.}
Sia $B = \texttt{MaxEnergyJoule}$ il budget dichiarato dall'utente.
Tutte le varianti con stima energetica superiore a $B + \varepsilon$,
dove $\varepsilon = 0.05 \cdot B$, vengono scartate a priori.
Le rimanenti vengono ripartite in due insiemi disgiunti:

\begin{equation}
  F_{\text{strict}} = \{\, v : E(v) \leq B \,\}
  \label{eq:f_strict}
\end{equation}
\begin{equation}
  F_{\text{relaxed}} = \{\, v : B < E(v) \leq B + \varepsilon \,\}
  \label{eq:f_relaxed}
\end{equation}

$F_{\text{strict}}$ contiene le varianti che rispettano il budget in modo
esatto; $F_{\text{relaxed}}$ quelle che lo superano entro una fascia del
5\%, ammesse come zona di tolleranza. Se entrambi gli insiemi sono vuoti,
l'invocazione viene rifiutata con HTTP~422.

\medskip
\noindent\textbf{Fase 2 — Identificazione dei candidati finali.}
Da ciascun insieme viene estratta la variante con errore minimo:

\begin{align}
  v_1 &= \arg\min_{v \in F_{\text{strict}}}\; \bigl(\mathit{Err}(v),\; E(v)\bigr)
  \label{eq:v1} \\[4pt]
  v_2 &= \arg\min_{v \in F_{\text{relaxed}}}\; \bigl(\mathit{Err}(v),\; E(v)\bigr)
  \label{eq:v2}
\end{align}

dove $\mathit{Err}(v)$ è il punteggio di errore della variante, definito
nella Sezione~\ref{sec:scheduler:error}, e il consumo energetico agisce
solo a parità di errore. Se uno dei due insiemi è vuoto, il
candidato corrispondente non esiste e la selezione finale è immediata.

\medskip
\noindent\textbf{Fase 3 — Confronto e selezione finale.}
Quando entrambi i candidati esistono, si pone la domanda: conviene
sforare il budget scegliendo $v_2$? La risposta è affermativa solo se il
guadagno relativo in termini di riduzione dell'errore compensa proporzionalmente
lo sforamento relativo rispetto al budget. La condizione è:

\begin{equation}
  \frac{\mathit{Err}(v_1) - \mathit{Err}(v_2)}{\mathit{Err}(v_1)}
  \;>\;
  \frac{E(v_2) - B}{B}
  \label{eq:tradeoff}
\end{equation}

Il lato sinistro esprime la riduzione relativa dell'errore ottenuta
passando da $v_1$ a $v_2$, normalizzata rispetto all'errore di $v_1$
stesso — il punto di riferimento naturale, dato che è la migliore
alternativa entro vincolo. Il lato destro esprime la percentuale di
sforamento energetico rispetto al budget dichiarato, valore compreso per
costruzione nell'intervallo $(0,\, 0.05]$.
%L'interpretazione è diretta: se $v_2$ sfora il budget del 3\%, è
%accettabile solo se riduce l'errore di oltre il 3\% rispetto a $v_1$.
%Una riduzione inferiore non giustifica la violazione del vincolo.
Se la condizione~\eqref{eq:tradeoff} è soddisfatta, si sceglie $v_2$; altrimenti,
 si sceglie $v_1$.

Si osservi che, quando $\mathit{Err}(v_1) = 0$, la scelta fra $v_1$ e $v_2$ è banalmente $v_1$
dato che rispetta il vincolo ed ha errore minimo.

\medskip
\noindent\textbf{Fase 4 — Gestione del fallimento.}
Se entrambi gli insiemi $F_{\text{strict}}$ e $F_{\text{relaxed}}$ sono
vuoti — ovvero tutte le varianti superano $B + \varepsilon$ —
l'invocazione viene rifiutata con HTTP~422, senza eseguire alcuna
computazione.

La logica complessiva delle due politiche è formalizzata
nell'Algoritmo~\ref{alg:select}.\\
L'implementazione dell'algoritmo riflette la scelta di mantenere separata
la logica di selezione della variante dalla logica di assegnazione del
container e di esecuzione: \textsc{SelectEnergyAwareVariant} restituisce
esclusivamente la variante selezionata, lasciando invariato il percorso di
esecuzione successivo.

Complessivamente, lo scheduler energy-aware introduce un livello di controllo sul trade-off tra consumo energetico e qualità del risultato che opera esclusivamente su base consensuale, rispetta i profili energetici aggiornati a runtime e si integra con il meccanismo di condivisione dei container tra varianti. Il risultato è un sistema adattivo che, pur rimanendo trasparente agli utenti che non desiderano intervenire sulla selezione, è in grado di soddisfare vincoli energetici espliciti minimizzando l'errore del risultato compatibilmente con i vincoli imposti e, a parità di accuratezza, preferendo sempre la variante meno costosa energeticamente.

\begin{algorithm}[H]
\caption{\texttt{SelectEnergyAwareVariant}: politica a casi}
\label{alg:select}
\begin{algorithmic}[1]
\Function{SelectEnergyAwareVariant}{$r$ : \texttt{*function.Request}}
    \State $\textit{base} \gets r.\texttt{Fun}$
    \State $V \gets \Call{GetFunctionsByLogicalName}{\textit{base}.\texttt{LogicalName}}$
    \If{$V = \emptyset$}
        \State \Return $(\textit{base},\; \text{``fallback-base''},\; \textbf{nil})$
        \Comment{nessuna variante registrata}
    \EndIf
    \State $\textit{Eval} \gets [\,]$
    \For{$fn \in V$}
        \State $\textit{warm} \gets \Call{HasWarmContainer}{fn}$
        \State $E(fn) \gets \Call{energyCostJoule}{fn,\; \textit{warm}}$
        \State $\mathit{Err}(fn) \gets \Call{errorScore}{fn}$
        \State $\textit{Eval}.\textsc{Append}(fn,\; E(fn),\; \mathit{Err}(fn))$
    \EndFor
    \If{$r.\texttt{MaxEnergyJoule} = \textbf{nil}$ \textbf{or} $r.\texttt{MaxEnergyJoule} \leq 0$}
        \Comment{Case: error-first}
        \State $v^* \gets \arg\min_{\textit{Eval}} (\mathit{Err},\, E)$
        \State \Return $(v^*,\; \text{``error-first''},\; \textbf{nil})$
    \EndIf
    \State $B \gets r.\texttt{MaxEnergyJoule}$
    \Comment{Caso B: energy-budget}
    \State $\varepsilon \gets \Call{energyEpsilon}{B}$
    \Comment{$\varepsilon = 0.05 \cdot B$}
    \State $F_{\text{strict}}   \gets \{x \in \textit{Eval} \mid x.E \leq B\}$
    \State $F_{\text{relaxed}} \gets \{x \in \textit{Eval} \mid B < x.E \leq B + \varepsilon\}$
    \If{$F_{\text{strict}} = \emptyset$ \textbf{and} $F_{\text{relaxed}} = \emptyset$}
        \State \Return $(\textbf{nil},\; \text{``no-feasible-variants''},\; \textsc{Error})$
        \Comment{HTTP 422}
    \EndIf
    \If{$F_{\text{strict}} = \emptyset$}
        \State $v^* \gets \arg\min_{F_{\text{relaxed}}} (\mathit{Err},\, E)$
        \State \Return $(v^*,\; \text{``relaxed-only''},\; \textbf{nil})$
    \EndIf
    \If{$F_{\text{relaxed}} = \emptyset$}
        \State $v^* \gets \arg\min_{F_{\text{strict}}} (\mathit{Err},\, E)$
        \State \Return $(v^*,\; \text{``strict-only''},\; \textbf{nil})$
    \EndIf
    \State $v_1 \gets \arg\min_{F_{\text{strict}}} (\mathit{Err},\, E)$
    \State $v_2 \gets \arg\min_{F_{\text{relaxed}}} (\mathit{Err},\, E)$
    \If{$\dfrac{\mathit{Err}(v_1) - \mathit{Err}(v_2)}{\mathit{Err}(v_1)} \geq \dfrac{E(v_2) - B}{B}$}
        \State \Return $(v_2,\; \text{``relaxed-tradeoff''},\; \textbf{nil})$
    \Else
        \State \Return $(v_1,\; \text{``strict''},\; \textbf{nil})$
    \EndIf
\EndFunction
\end{algorithmic}
\end{algorithm}

